%% kapa: exam=1 round=30 year=2019 period=2 form=A subject=appraisal_law subject_ko=감정평가관계법규 seq=29 number=29
%% source: Q-Net 감정평가사 기출문제 (cst00309) artlSeq=5205205 file=extracted/1008/5205205/2019년 제30회 감정평가사 1차시험 문제지/2019년 제30회 감정평가사 1차시험 문제지/2019년도 제30회 감정평가사 1차 2교시 A형.hwp
%% conversion: kordoc 4.24.51 -> tools/gen1.py (auto); review=proofread-sonnet; flags=none
\begin{question}{29}
공간정보의 구축 및 관리 등에 관한 법령상 지적공부에 관한 설명으로 옳지 않은 것은?
\begin{choices}
  \item 정보처리시스템을 통하여 기록ㆍ저장한 지적공부의 전부가 멸실된 경우에는 국토교통부장관은 지체 없이 이를 복구하여야 한다.
  \item 국토교통부장관은 정보처리시스템을 통하여 지적정보관리체계에 기록ㆍ저장한 지적공부가 멸실될 경우를 대비하여 지적공부를 복제하여 관리하는 정보관리체계를 구축하여야 한다.
  \item 지적공부를 정보처리시스템을 통하여 기록ㆍ저장한 경우 관할 시ㆍ도지사, 시장ㆍ군수 또는 구청장은 그 지적공부를 지적정보관리체계에 영구히 보존하여야 한다.
  \item 국토교통부장관은 지적공부의 효율적인 관리 및 활용을 위하여 지적정보 전담 관리기구를 설치ㆍ운영한다.
  \item 지방자치단체의 장이 지적전산자료를 신청하는 경우에는 지적전산자료의 이용 목적 등에 관하여 미리 관계 중앙행정기관의 심사를 받지 않아도 된다.
\end{choices}
\end{question}
\answer{1}
